\documentclass[10pt,a4paper]{amsart}
\usepackage[margin=19mm]{geometry}
\usepackage{amsmath,amssymb,mathtools}
\usepackage[scr=rsfs,cal=euler]{mathalfa}
\usepackage{xfrac}
\usepackage[unicode,hidelinks,bookmarks=false]{hyperref}
\newcommand{\ie}{i.e.\ }
\newcommand{\cf}{cf.\ }
\newcommand{\resp}{respectively\ }
\newcommand{\Subsection}{\subsection}
\providecommand{\nobreakdash}{\nobreak\hbox{-}}
\setlength{\emergencystretch}{2em}
\multlinegap0pt
\usepackage{bm}
\usepackage{bbm}
\usepackage{dsfont}
\usepackage{xspace}
\usepackage{cancel}
\usepackage{mathtools}
\usepackage{amsthm}
\makeatletter
\@ifundefined{proposition}{\newtheorem{proposition}{Proposition}}{}
\makeatother
\newcommand{\Loss}{\mathcal{L}}
\title[Do Deep Networks Forget Initialization? A Forgetting-Time Vi]{Do Deep Networks Forget Initialization? A Forgetting-Time View of Practical Inductive Bias}
\author{Mohua Das, Pierfrancesco Beneventano, Shibshankar Dey, Gareth H. McKinkey, Tomaso Poggio}
\date{Source: arXiv:2605.29152v1 (CC BY 4.0)}
\begin{document}
\maketitle

\noindent\emph{Source and licence.} Do Deep Networks Forget Initialization? A Forgetting-Time View of Practical Inductive Bias. Version arXiv:2605.29152v1 (CC BY 4.0), distributed under the Creative Commons Attribution 4.0 International licence, as recorded in the arXiv licence metadata for this exact version and shown on the abstract page https://arxiv.org/abs/2605.29152v1. Full rights notice: licenses/arxiv-2605.29152-CC-BY-4.0.txt.\par\smallskip

\noindent\textit{Editorial scope.} Counted unit: the Proposition whose source label is \texttt{None} in the article (source0.tex lines 2397--2420, source label \texttt{None}), delivered together with the article's own complete original proof of it (source0.tex lines 2422--2533, one \texttt{proof} environment of 112 lines). No mathematics is written, repaired or re-derived: statement and proof are the article's own text line for line. Required context retained uncounted: none. Every definition, display and label of those lines is retained unchanged. Omitted from this package and not redistributed: 1--2396, 2421--2421, 2534--3146. Nothing inside a retained excerpt is omitted or altered: every retained body is delivered exactly as the article's lines are written, so no omission receipt of any kind accompanies this package. Citations: the retained excerpts carry no citation command at all, so no bibliography entry of the article is transcribed and no reference is redistributed. Reference adaptation: none was needed and none was made; every reference inside the delivered unit resolves inside it, so no unresolved reference or citation is produced. Printed slips kept literal: no printed word, symbol, hypothesis or number of the counted statement or of its proof was altered. Layout and class: the article's own class is \texttt{article} with packages including neurips\_2026, inputenc, fontenc, hyperref, url, booktabs, amsmath, amssymb, amsfonts, mathtools, amsthm, microtype, xcolor, graphicx; the wrapper is the lane's pinned \texttt{amsart} editorial wrapper at 10pt on A4, which restates the theorem-like environments the retained text needs and adds the article's own macro definitions that the retained text needs (\texttt{\textbackslash Loss}), with extra wrapper packages (bm, bbm, dsfont, xspace, cancel, mathtools). The delivered numbering is the unit's own. Assets: no retained span contains a figure environment, a graphics-inclusion command, a PDF-page inclusion, a table or a TikZ picture; no image file is copied into this package, the package declares no asset dependency and no image file is redistributed with it. Licence: version arXiv:2605.29152v1 of this article is distributed under the Creative Commons Attribution 4.0 International licence, as recorded in the arXiv licence metadata for this exact version and shown on the abstract page https://arxiv.org/abs/2605.29152v1. A full rights notice is delivered at licenses/arxiv-2605.29152-CC-BY-4.0.txt. Characters: every retained excerpt is pure ASCII, so the delivered engine, XeLaTeX with its default Latin Modern text font, reports no missing character.
\begin{proposition}[Deep linear gradient flow preserves adjacent imbalances]
Assume \(\phi\) is differentiable.  Along Euclidean gradient flow
\[
    \dot W_j=-\nabla_{W_j}\Loss,
    \qquad
    j=1,\ldots,L,
\]
one has
\[
    \frac{d}{dt}D_j(t)=0,
    \qquad
    j=1,\ldots,L-1.
\]
With coupled continuous-time weight decay,
\[
    \dot W_j=-\nabla_{W_j}\Loss-\lambda W_j,
\]
the imbalances satisfy
\[
    \frac{d}{dt}D_j(t)=-2\lambda D_j(t),
    \qquad
    D_j(t)=e^{-2\lambda t}D_j(0).
\]
\end{proposition}

\begin{proof}
For each layer \(j\), define
\[
    A_j=W_LW_{L-1}\cdots W_{j+1},
    \qquad
    B_j=W_{j-1}\cdots W_1,
\]
with the convention that empty products are identity matrices.  Then
\[
    F=A_jW_jB_j.
\]
Let
\[
    G=\nabla_F\phi(F).
\]
By the chain rule,
\[
    \nabla_{W_j}\Loss
    =
    A_j^\top G B_j^\top.
\]
We use the identities
\[
    A_j=A_{j+1}W_{j+1},
    \qquad
    B_{j+1}=W_jB_j.
\]
Let
\[
    H_j=\nabla_{W_j}\Loss=A_j^\top G B_j^\top.
\]
Then
\[
\begin{aligned}
    H_jW_j^\top
    &=
    A_j^\top G B_j^\top W_j^\top  \\
    &=
    W_{j+1}^\top A_{j+1}^\top G B_{j+1}^\top  \\
    &=
    W_{j+1}^\top H_{j+1},
\end{aligned}
\]
and
\[
\begin{aligned}
    W_jH_j^\top
    &=
    W_jB_jG^\top A_j  \\
    &=
    B_{j+1}G^\top A_{j+1}W_{j+1}  \\
    &=
    H_{j+1}^\top W_{j+1}.
\end{aligned}
\]
Therefore, under gradient flow,
\[
\begin{aligned}
    \frac{d}{dt}(W_jW_j^\top)
    &=
    -H_jW_j^\top-W_jH_j^\top  \\
    &=
    -W_{j+1}^\top H_{j+1}
    -
    H_{j+1}^\top W_{j+1}.
\end{aligned}
\]
On the other hand,
\[
\begin{aligned}
    \frac{d}{dt}(W_{j+1}^\top W_{j+1})
    &=
    -H_{j+1}^\top W_{j+1}
    -
    W_{j+1}^\top H_{j+1}.
\end{aligned}
\]
The two derivatives are identical, so
\[
    \frac{d}{dt}
    \left(
        W_jW_j^\top
        -
        W_{j+1}^\top W_{j+1}
    \right)
    =
    0.
\]
This proves conservation.

With weight decay, the gradient-dependent terms are unchanged and still
cancel.  The additional terms are
\[
    -2\lambda W_jW_j^\top
    \qquad\text{and}\qquad
    -2\lambda W_{j+1}^\top W_{j+1},
\]
so
\[
    \dot D_j
    =
    -2\lambda
    \left(
        W_jW_j^\top
        -
        W_{j+1}^\top W_{j+1}
    \right)
    =
    -2\lambda D_j.
\]
Solving gives \(D_j(t)=e^{-2\lambda t}D_j(0)\).
\end{proof}

\end{document}
